\lysh{35479}{4}

\begin{alist}
\item Θα πρέπει να ισχύει: $F_4 = 1 \iff \frac{25}{\lambda} = 1 \iff \lambda = 25$.

\item Έχουμε: $F_1 = \frac{4}{25} = f_1$.

$F_2 = \frac{11}{25}$, αλλά $F_2 = F_1 + f_2 \iff \frac{11}{25} = \frac{4}{25} + f_2 \iff f_2 = \frac{7}{25}$.

$F_3 = \frac{18}{25}$, αλλά $F_3 = F_2 + f_3 \iff \frac{18}{25} = \frac{11}{25} + f_3 \iff f_3 = \frac{7}{25}$.

$F_4 = F_3 + f_4 \iff 1 = \frac{18}{25} + f_4 \iff f_4 = \frac{7}{25}$.

\item για να υπολογίσουμε τη μέση τιμή, θα χρησιμοποιήσουμε τον τύπο:
\begin{align*}
\bar{x} &= x_1 \cdot f_1 + x_2 \cdot f_2 + x_3 \cdot f_3 + x_4 \cdot f_4 = \\
&= \alpha \cdot \frac{4}{25} + (\alpha + 5) \cdot \frac{7}{25} + (\alpha + 10) \cdot \frac{7}{25} + (\alpha + 35) \cdot \frac{7}{25} = \\
&= \frac{4\alpha}{25} + \frac{7\alpha + 35}{25} + \frac{7\alpha + 70}{25} + \frac{7\alpha + 245}{25} = \frac{25\alpha + 350}{25} = \frac{25 \cdot (\alpha + 14)}{25} = \alpha + 14
\end{align*}
Αφού $\bar{x} = 19 \iff \alpha + 14 = 19 \iff \alpha = 5$.
\end{alist}
\vspace{6mm}
\noindent\rule{\textwidth}{0.4pt}
\vspace{6mm}

